\chapter{Coupled Neutronic--Thermal-Hydraulic Dynamics}
\label{ch:coupled}

We now assemble the pieces---point kinetics, lumped thermal model, and feedback
coefficients---into a single coupled dynamical system and explore its qualitative
behaviour before the formal stability analysis of Part~III.

\section{The closed-loop model}

Combining the point kinetics \eqref{eq:prk}, the lumped thermal model of
Chapter~\ref{ch:thermal}, and the feedback law \eqref{eq:rho-feedback}, a compact
representative model is
\begin{align}
\frac{\dd p}{\dd t} &= \frac{\rho(t)-\beta}{\Lambda}\,p + \sum_i \lambda_i \zeta_i,\\
\frac{\dd \zeta_i}{\dd t} &= \frac{\beta_i}{\Lambda}\,p - \lambda_i \zeta_i,\\
C_f\frac{\dd \Tf}{\dd t} &= P_0\,p - U\!A\,(\Tf-\Tc),\\
C_c\frac{\dd \Tc}{\dd t} &= U\!A\,(\Tf-\Tc) - 2\dot m c_c(\Tc-T_{in}),\\
\rho(t) &= \rho_{ext}(t) + \alpha_D(\Tf-\Tf^0) + \alpha_c(\Tc-\Tc^0).
\end{align}
This is a system of nonlinear ordinary differential equations: nonlinear because
reactivity multiplies power in the first equation, and reactivity itself depends
on the temperatures driven by power. The nonlinearity is essential---it is what
makes a power excursion turn over rather than grow forever---but it also makes the
behaviour rich.

\section{The self-regulating reactor in words}

Before any mathematics, the qualitative behaviour of a \emph{negative-feedback}
reactor is worth stating plainly. Suppose the operator withdraws a control rod,
inserting positive $\rho_{ext}$. Power rises. The fuel heats within milliseconds;
if $\alpha_D<0$, this immediately subtracts reactivity, blunting the rise. As the
coolant heats over the next second, $\alpha_c<0$ subtracts more. The reactor
settles at a higher power and temperature at which the feedback exactly cancels
the inserted reactivity: a new stable operating point. The reactor has
\emph{regulated itself} without any control action. Conversely, if the operator
reduces power, temperatures fall, feedback adds reactivity, and the fall is
arrested. This is the everyday meaning of a negative power coefficient.

Now reverse the sign of the dominant feedback, as in the RBMK at low power. The
rod withdrawal raises power; boiling increases; the positive void coefficient
adds \emph{more} reactivity; power rises further; and unless a strong prompt
negative effect intervenes, the excursion accelerates. The same control action
that a stable reactor absorbs, the unstable one amplifies.

\section{Equilibrium and the power defect}

At equilibrium the time derivatives vanish and the feedback exactly offsets the
external reactivity. The total reactivity that must be supplied by control to
hold the reactor critical from cold-zero-power to hot-full-power is the
\textbf{power defect},
\begin{equation}
\Delta\rho_{\text{defect}} = -\!\int_0^{P_0}\!\alpha_P\,\dd P,
\end{equation}
which for a negative power coefficient is positive (control must \emph{add}
reactivity to reach full power, because the feedback removes it). The power defect
is a useful integral measure of how strongly self-regulating a reactor is, and it
must be matched by available control-rod and chemical-shim worth.

\section{Linearisation preview}

For small perturbations about an operating point, write $p=p_0+\delta p$,
$\Tf=\Tf^0+\delta \Tf$, etc., and keep first-order terms. The linearised system
has the form $\dot{\mathbf x}=\mathbf A\,\mathbf x + \mathbf b\,\rho_{ext}$, and
its stability is decided by the eigenvalues of the matrix $\mathbf A$: the
operating point is stable if and only if all eigenvalues have negative real part.
The feedback coefficients enter $\mathbf A$ directly, so the sign and magnitude of
$\alpha_D$, $\alpha_c$, and $\alpha_v$ determine the eigenvalues' signs. This is
the formal content of Chapter~\ref{ch:linear}, where we also pass to the frequency
domain and the Nyquist criterion, which handles the transport delays more
naturally than eigenvalues alone.

\section{A demonstration of self-limitation}

Figure~\ref{fig:transients} (bottom right) already showed the canonical result:
a reactivity-initiated excursion that, with negative feedback, rises to a peak and
then turns over as the heating fuel removes reactivity. The peak power, the time
to peak, and the total energy released are set by the competition between the
inserted reactivity and the feedback strength---precisely the Nordheim--Fuchs
balance derived in Chapter~\ref{ch:nonlinear}. The same model with the feedback
sign flipped does not turn over; it diverges until the model's assumptions break.
The entire safety case of an inherently stable reactor rests on getting that sign
right and the magnitude large enough.

\keypoint{The coupled model is nonlinear, but its stability is governed by a
simple principle: the feedback must remove reactivity as power and temperature
rise, and it must do so promptly enough. Part~III makes ``promptly enough''
precise through eigenvalues and the Nyquist criterion; Part~IV shows what the
violation of this principle looked like in practice.}


\section*{Exercises}
\addcontentsline{toc}{section}{Exercises}
\begin{enumerate}[leftmargin=*,itemsep=2pt]
\item Describe in words the self-regulating response of a negative-feedback reactor to a rod withdrawal.
\item Define the power defect and explain its sign for a negative power coefficient.
\item Explain how the feedback coefficients enter the linearised system matrix and why feedback is more potent at higher power.
\item Sketch the qualitative power response of a reactor with positive feedback to a small rod withdrawal.
\item Why does a nonlinear feedback model produce a power excursion that turns over, while a linear analysis cannot capture the turnover?
\end{enumerate}
